\clearpage\SourceLeaf{302}
\begin{center}
\textbf{[ZU MÖBIUS' BARYCENTRISCHEM CALCUL.]}

\medskip

[I.]
\end{center}

\begin{center}\textsc{Gauss} an \textsc{Schumacher}. Göttingen, 12. Mai 1843.\end{center}

Ich bin Ihnen, mein theuerster Freund, noch für mehrere Ihrer Briefe und sonstige
interessante Mittheilungen meinen Dank schuldig. Sie dürfen jedoch nicht zu genau mit
mir rechnen: seit 4---6 Wochen bin ich (anfangs durch zufällige Umstände) in einige
mathematische Speculationen hineingezogen, wo ich immer wieder durch neue Aussichten in
andere Richtungen gelenkt wurde, und vieles erreicht, vieles verfehlt habe. In einem
solchen Herumtreiben ist man (oder bin ich wenigstens) fast unfähig, mich in irgend
andere Gegenstände einzulassen, und so ist namentlich auch in meiner Correspondenz ungemein viel Rückstand geworden.
Jene Speculationen betrafen grossentheils weniger neue Sachen, als Durchführung mir
eigenthümlicher Methoden; zuletzt u. a. mehreres sich auf die Kegelschnitte
beziehendes. Mir ist dabei wieder in Erinnerung gekommen, dass ich vor einem halben
Jahrhundert, als ich zuerst \textsc{Newtons} Principia las, mehreres unbefriedigend
fand, namentlich auch seine an sich herrlichen Sätze die Kegelschnitte betreffend.
Aber ich las immer mit dem Gefühle, dass ich durch das Erlernte nicht Herr der Sache
wurde; besonders quälte mich die gerade Linie, mit deren Hülfe ein Kegelschnitt
beschrieben werden kann. \textsc{Newton} löst die Aufgabe: durch 5 Punkte $A$, $B$,
$C$, $D$, $E$ einen Kegelschnitt zu legen, so auf, dass er erst vermittelst der vier
